\chapter{Head and Neck}

\label{ch:ch39_head_neck}

The head and neck contain many complex structures packed into a small space---the eyes, ears, sinuses, throat, salivary glands, thyroid, and major blood vessels. Imaging helps doctors evaluate problems such as hearing loss, sinus infections, neck lumps, sore throat that does not resolve, and tumours. The main imaging tools are CT (computed tomography), MRI (magnetic resonance imaging), and ultrasound.

\textbf{CT scans} of the head and neck are quick and excellent for looking at bones (such as the tiny bones of the ear) and detecting calcium in tumours. \textbf{MRI scans} provide superior detail of soft tissues---the brain, nerves, muscles, and tumours---and are preferred when the question involves the eye sockets, the base of the skull, or the extent of a known tumour. \textbf{Ultrasound} is useful for evaluating neck lumps, thyroid nodules, and salivary glands without radiation. In some cases, contrast dye is injected to highlight blood vessels or tumours.

Your doctor will choose the right test based on your symptoms. Always inform the imaging team about any metal implants, pacemakers, or if you might be pregnant.

---

\section{Coverage}
Head and neck imaging may include face, orbit, temporal bone, sinonasal region, pharynx, larynx, neck spaces, thyroid, salivary glands, and cervical vessels. Protocol selection depends on the clinical question and whether bone, mucosa, soft tissue, vascular anatomy, or nerves are the target.

\begin{table}[h]
\centering
\caption{Modality selection by anatomical region}
\begin{tabular}{p{3.5cm} p{3.5cm} p{3.5cm} p{3.5cm}}
\hline
\textbf{Region} & \textbf{CT} & \textbf{MRI} & \textbf{Ultrasound} \\
\hline
Carotid / vertebral & CTA (stenosis, dissection) & MRA (no radiation) & Doppler (stenosis screening) \\
Larynx & Cartilage invasion & Vocal cord paralysis, tumour depth & N/A \\
Neck masses & Staging, deep extension & Nerve sheath vs lymph node & First-line (superficial) \\
Oral cavity / oropharynx & Staging (bone invasion) & Perineural spread, deep tongue base & N/A \\
Orbits & Fractures, foreign bodies & Tumour extent, optic neuritis, cavernous sinus & Thyroid orbitopathy \\
Salivary glands & Calculi (non-contrast CT) & Tumour characterization & First-line (sialadenitis, masses) \\
Sinuses & First-line (acute sinusitis, polyps) & Complications, fungal sinusitis & N/A \\
Temporal bone & First-line (fractures, cholesteatoma) & Limited role & N/A \\
Thyroid & Not first-line & Limited & First-line (nodules, FNA guidance) \\
\hline
\end{tabular}
\end{table}

---

\section{Temporal Bone Imaging}
CT of the temporal bone is performed with thin cuts ($\leq 0.625$ mm) through the petrous temporal bone. Key structures include:

\begin{itemize}
  \item \textbf{External auditory canal}: Exostoses, canal cholesteatoma
  \item \textbf{Tympanic membrane and middle ear}: Ossicular chain integrity, cholesteatoma (soft tissue mass with ossicular erosion), tympanosclerosis
  \item \textbf{Inner ear}: Cochlear malformations, semicircular canal dehiscence, labyrinthine ossification
  \item \textbf{Internal auditory canal}: Widening suggests vestibular schwannoma (compare with MRI for soft tissue characterization)
  \item \textbf{Facial nerve canal}: Dehiscence or erosion in infection/tumour
\end{itemize}

\begin{pearlbox}[Attic cholesteatoma on CT]
A soft tissue mass in the epitympanum (attic) with erosion of the scutum and ossicular chain is highly suggestive of cholesteatoma. CT defines the extent of bone destruction; MRI (particularly diffusion-weighted imaging with non-EPI sequences) confirms the diagnosis and differentiates cholesteatoma from granulation tissue.
\end{pearlbox}

---

\section{Orbital Imaging}
Orbital emergencies include \textbf{orbital compartment syndrome} (proptosis, decreased vision, afferent pupillary defect from retrobulbar hemorrhage) and \textbf{orbital cellulitis} (proptosis, extraocular muscle enlargement, periorbital fat stranding, potential abscess).

For tumours, MRI is preferred to delineate the relationship to the optic nerve, extraocular muscles, and intracranial extension. Common orbital masses include cavernous haemangioma (well-defined, progressive enhancement), optic nerve glioma, optic nerve sheath meningioma (tram-track sign on axial views), and lacrimal gland tumours.

---

\section{Sinonasal Imaging}
CT is the first-line investigation for sinus disease. Review the ostiomeatal unit (OMU) for obstruction that predisposes to recurrent sinusitis.

\begin{table}[h]
\centering
\caption{Sinonasal CT findings}
\begin{tabular}{p{4cm} p{4.5cm} p{4.5cm}}
\hline
\textbf{Finding} & \textbf{Significance} & \textbf{Next step} \\
\hline
Aggressive bone destruction & Malignancy (e.g.\ squamous cell carcinoma) & Biopsy, staging MRI \\
Air--fluid levels & Acute sinusitis & Antibiotics, drainage if refractory \\
Fungal debris (hyperdense on CT) & Fungal sinusitis & Surgical debridement, biopsy \\
Mucosal thickening & Chronic rhinosinusitis & Medical management \\
Opacification with bone expansion & Mucocele & Surgical drainage \\
\hline
\end{tabular}
\end{table}

---

\section{Neck Masses and Staging of Head and Neck Cancer}
A systematic approach to a neck mass includes:

\begin{enumerate}
  \item \textbf{Location}: Which neck level (Levels I--VII)?
  \item \textbf{Laterality}: Unilateral vs bilateral
  \item \textbf{Characteristics}: Cystic vs solid, enhancing, necrotic, calcified
  \item \textbf{Relationship to adjacent structures}: Vessel displacement, nerve involvement, bone invasion
\end{enumerate}

For known or suspected head and neck squamous cell carcinoma, staging requires:
\begin{itemize}
  \item \textbf{MRI of the primary site}: Best for evaluating depth of invasion, perineural spread, and marrow involvement
  \item \textbf{CT of the neck with contrast}: Lymph node staging (size, necrosis, extracapsular spread)
  \item \textbf{CT chest}: Exclude pulmonary metastases
  \item \textbf{PET-CT}: For advanced disease (stage III--IV) to identify distant metastases
\end{itemize}

\begin{pearlbox}[Necrotic lymph node]
A lymph node with central low attenuation (necrosis), irregular enhancement, and loss of the fatty hilum is suspicious for metastatic squamous cell carcinoma regardless of size. Extracapsular spread (stranding of the surrounding fat) upstages the disease and may warrant adjuvant chemoradiotherapy.
\end{pearlbox}

---

\section{Laryngeal Imaging}
CT is the primary modality for laryngeal cancer staging. Key features to report:
\begin{itemize}
  \item \textbf{Tumour site}: Supraglottic, glottic, or subglottic
  \item \textbf{Cartilage invasion}: Erosion or marrow signal change on MRI; sclerosis or lysis on CT
  \item \textbf{Vocal cord fixation}: Correlates with deep muscle invasion
  \item \textbf{Pre-epiglottic / paraglottic space invasion}: Upstages the tumour
  \item \textbf{Airway compromise}: May require urgent intervention
\end{itemize}

---

\section{Interpretation}
Review compartments, mucosal surfaces, lymph-node stations, airway, vessels, bones, teeth, and skull base. Describe laterality, deep-space involvement, perineural spread, and relationship to critical structures when relevant.

\section{Safety}
Consider dental artifact, contrast, renal function, pregnancy, airway risk, and MRI implants. Suspicious airway compromise, vascular injury, orbital compartment syndrome, or aggressive infection requires urgent escalation.

\clearpage
